% Teil II: durchgerechneter Weg, Teile a--c.

\subsection*{Gegeben und gesucht}
Zwei gleiche Pendel ($m_1=m_2=m$, $D_1=D_2=D$), gekoppelt über Feder $k$.
Gesucht: (a) Bewegungsgleichungen, (b) Normalschwingungen, (c) $x_1(t),x_2(t)$
für die Anfangsbedingung $x_1(0)=x_{1,0}$, $x_2(0)=\dot x_1(0)=\dot x_2(0)=0$
und die Energietransfer-Frequenz.

\subsection*{Rechnung}
\textbf{a) Bewegungsgleichungen.} Kräfte: eigene Feder $D\,x_i$ (rücktreibend)
und Koppelfeder $k(x_2-x_1)$:
\begin{equation}
  m\,\ddot x_1=-D\,x_1+k(x_2-x_1),\qquad
  m\,\ddot x_2=-D\,x_2-k(x_2-x_1)= -D\,x_2+k(x_1-x_2)\;.
\end{equation}

\textbf{b) Entkopplung und Normalschwingungen.} Addieren der beiden
Gleichungen: $m(\ddot x_1+\ddot x_2)=-D(x_1+x_2)$. Subtrahieren:
$m(\ddot x_1-\ddot x_2)=-D(x_1-x_2)+2k(x_2-x_1)=-(D+2k)(x_1-x_2)$. Mit
$y_\pm=\tfrac12(x_1\pm x_2)$:
\begin{equation}
  \ddot y_++\omega_+^2y_+=0,\qquad \ddot y_--+\omega_-^2y_-=0,
  \quad
  \boxed{\;\omega_+=\sqrt{\frac{D}{m}}\;},\quad
  \boxed{\;\omega_-=\sqrt{\frac{D+2k}{m}}\;}.
\end{equation}
Allgemein: $y_\pm(t)=A_\pm\cos(\omega_\pm t+\varphi_\pm)$
(gleichphasiger Oszillator $\omega_+$: Koppelfeder bleibt ungedehnt;
gegenphasiger Oszillator $\omega_-$: beide Koppelfedern wirken).

\textbf{c) Anfangsbedingungen.} $y_+(0)=\tfrac{x_{1,0}}2$, $y_-(0)=\tfrac{x_{1,0}}2$,
$\dot y_+(0)=\dot y_-(0)=0$, also $A_\pm=x_{1,0}/2$, $\varphi_\pm=0$:
\begin{equation}
  y_\pm(t)=\frac{x_{1,0}}2\cos(\omega_\pm t)\;.
\end{equation}
Rücktransformation mit $\omega_{\mathrm m}=\frac{\omega_-+\omega_+}2$ und
$\Delta=\frac{\omega_--\omega_+}2$:
\begin{align}
  x_1(t)&=y_++y_-=\frac{x_{1,0}}2\big[\cos\omega_+t+\cos\omega_-t\big]
         =\boxed{\;x_{1,0}\,\cos(\omega_{\mathrm m}t)\,\cos(\Delta t)\;},\\
  x_2(t)&=y_+-y_-=\frac{x_{1,0}}2\big[\cos\omega_+t-\cos\omega_-t\big]
         =\boxed{\;-x_{1,0}\,\sin(\omega_{\mathrm m}t)\,\sin(\Delta t)\;}.
\end{align}

\textbf{Energietransfer.} Die Amplitude von $x_1$ moduliert mit
$|\cos\Delta t|$, die zugehörige Energie wie $\cos^2\Delta t$; die Energie
schwappt also periodisch mit
\begin{equation}
  \boxed{\;\omega_{\text{Energie}}=2\Delta=\omega_- -\omega_+\;}
  \qquad\text{bzw.}\quad
  f_{\text{Energie}}=\frac{\omega_--\omega_+}{2\pi}
  =\Big(\sqrt{\tfrac{D+2k}{m}}-\sqrt{\tfrac{D}{m}}\Big)\frac1{2\pi}\;.
\end{equation}
Anfang: ganze Energie in Masse 1; nach $T=\pi/\Delta=2\pi/(\omega_--\omega_+)$
ist sie vollständig zur Masse 2 gewandert und kommt mit derselben Periode
zurück.

\begin{ergebnis}
$\omega_+=\sqrt{D/m}$, $\omega_-=\sqrt{(D+2k)/m}$;
$x_1=x_{1,0}\cos(\omega_{\mathrm m}t)\cos(\Delta t)$,
$x_2=-x_{1,0}\sin(\omega_{\mathrm m}t)\sin(\Delta t)$;
Energietransfer-Frequenz $=(\omega_--\omega_+)/2\pi$.
\end{ergebnis}

\begin{kontrolle}
\textbf{Grenzfall $k=0$:} $\omega_+=\omega_-=\sqrt{D/m}$, die Massen schwingen
ungestört; $x_1=x_{1,0}\cos\omega_0t$, $x_2=0$ --- die Anfangsbedingung.
\textbf{Querprobe:} $x_1(0)=x_{1,0}$ und $x_2(0)=0$ erfüllt; $\dot x_1(0)$ und
$\dot x_2(0)$ verschwinden (Kosinus-/Sinus-Anfang). \checkmark

\textbf{Einheiten:} $\omega_\pm$ in \unit{\per\second}, Frequenz in
\unit{\hertz}. \checkmark
\end{kontrolle}

\begin{merksatz}
Gekoppelte Oszillatoren löst man über Summe/Differenz (Normalkoordinaten). Die
Schwebung der Einzelmassen und der hin- und herwandernde Energieanteil
laufen mit der \emph{Differenz} der beiden Normalfrequenzen $\omega_--\omega_+$.
\end{merksatz}
